% !TEX root = ../main.tex
\pagebreak
\lecture{7}{2026-10-29}{Marketingstrategie}
\textit{siehe Folien} \url{https://lms.uzh.ch/auth/RepositoryEntry/17910989063/CourseNode/96531468745796/path%3D~~Vorlesung%202/0}
\subsection{Marketingstrategie Vorlesung}

Strategieentwicklung umfasst 2 Fragen
\begin{enumerate}
    \item Zielmarkt: Welche Kunden wollen wir (nicht) bedienen?
    \item Wie positionieren sich unternehmen? Worin liegt Nutzenversprechen?
\end{enumerate}

\subsection{Vorlesung: Was ist eine Haltung wert (in \$)?}
Fallstudie: Verträge KI unternehmen und ''Department of War'', worin Anthropic und OpenAI den Einsatz ihrer KI-Modelle in autonomen Waffen und Massenüberwachung ablehnten. OpenAI unterschrieb jedoch einen angepassten Vertrag, während Anthropic kein Vertrag abschloss;

\textbf{Konsequenzen für unternehmen:}
\begin{itemize}
    \item Anthropic: verlor Vertrag im Wert hunderte Mio.
    \item OpenAI: wurde von Konsumenten boykottiert
\end{itemize}

\subsubsection{Mission/Vision Statements}
\begin{itemize}
    \item Anthropic: ''AI Safety and research company. We build reliable,...,steerable AI systems''
    \begin{itemize}
        \item Sehr realistisch, praxisorientiert. Das Wie steht im Vordergrund.
    \end{itemize}
    \item ''OpenAI: ensure that general AI benefits all of humanity'' \footnote{Die Unterschreibung des militär Vertrags widerspricht die Mission und Werte}
    \begin{itemize}
        \item Beschreibt nur ein ideales Bild nicht wie es erreicht wird
    \end{itemize}
\end{itemize}

Anthropic monetarisiert wenige, tiefe Beziehungen und OpenAI oberflächliche Reichweite.


\subsubsection{Ursprung des Konflikts}
Aufgrund der Verkaufszahlen von Anthropic könnten sich diese Positionierung leisten, weil alle ihre Nutzer bereits für die Nutzung zahlen. 
Da OpenAI meist kostenlos angeboten wird, schien ein solches Angebot sehr attraktiv auch relativ zum Verlust von Kunden. OpenAI hatte aufgrund des Erlösmodells weniger Entscheidungsmacht in Bezug auf die Gewinnmaximierung.

Nach 48h-Sperrung von Anthropic durch die Regierung, unterzeichnete OpenAI den Vertrag. Claude wurde zur Nr.1 App in Appstore (als moralisch korrekte Marke betrachtet) und OpenAI wird von der Öffentlichkeit boykottiert. OpenAI kommt aus  
Dabei ging es aber nicht rein um moralische Werte, sondern um Abwägung von Cost/Return.
\pagebreak
\subsubsection{Gleiche Werte, entgegengesetzte Entscheidungen}

\textbf{Anthropic}

\begin{itemize}
    \item Zwei rote Linien
    \item 2 Gründe fürs Ablehnen: demokratische Werte und Zuverlässigkeit
    \item Kosten: 950 Mio
\end{itemize}

Das Erlösmodell beschränkt die strategischen Entscheidungungsmöglichkeiten eines Unternehmens.

\subsubsection{Schwartz-Werte der Parteien}

\begin{itemize}
    \item Anthropic
\begin{itemize}
    \item Universalismus, Benevolenz, Selbstbestimmung
    \item Also eine zusammenarbeit mit dem DoD würde einen Widerspruch zu seinen Werten und die ihrer Kunden
\end{itemize}
\item OpenAI
\begin{itemize}
    \item Fast gleich wie Anthropic
\end{itemize}
\item Department of War
\begin{itemize}
    \item Macht, Konformität, Sicherheit
\end{itemize}
\end{itemize}

\textbf{Drei Thesen}

\begin{itemize}
    \item Konflikt strukturell anhand der kommunizierten Werte unlösbar
    \item OpenAI widersprach ihre Werte in dem sie beide Werte Pole besetzte
    \item Anthropic baute Brücke gezielt zum Gegenpol
    \begin{itemize}
        \item Das Demokratie-Argument bedient Universalismus. Das Argument «not reliable enough» bedient Sicherheit — einen Wert der konservativeren Gegenseite.
    \end{itemize}
\end{itemize}

\subsubsection{Debatte von studis nach Reihenfolge}

\textbf{Geschäftsleitung Anthropic}
\begin{itemize}
    \item Statement an Geschäftsleitung mit Fakten (Kennzahl) 
    \begin{itemize}
        \item Wir können Massenüberwachung nicht vertreten
        \item Geschäftliche Gründe: Response von Konsumenten war das Ziel
        \item Fakt: Nr.1 App in Appstore
    \end{itemize}
    \item konkrete Forderung an Geschäftsleitung
    \begin{itemize}
        \item keine Aussage
    \end{itemize}
    \item Was wäre das stärkste Argument der Gegenseite
    \begin{itemize}
        \item keine Aussage
    \end{itemize}
    \item Eine Erfolgskennzahl der letzten 12 Monate
    \begin{itemize}
        \item keine Aussage
    \end{itemize}
\end{itemize}
\textbf{Department of War (Politik)}
\begin{itemize}
    \item Statement an Geschäftsleitung mit Fakten (Kennzahl)
          \begin{itemize}
              \item Wir beauftragten 4 KI Anbieter bis 200 Mio mit demselben Vertrag
              \item Demokratische rechtsstaatliche Bestimmung, die nicht von einem einzigen Anbieter abzulehnen sei
              \item Man wirft uns vor Überwachung zu wollen, das stimmt nicht
          \end{itemize}
    \item konkrete Forderung an Geschäftsleitung
          \begin{itemize}
              \item Alle Anbieter sollten die Vertragsbedingungen akzeptieren 
          \end{itemize}
    \item Was wäre das stärkste Argument der Gegenseite
          \begin{itemize}
              \item
          \end{itemize}
    \item Eine Erfolgskennzahl der letzten 12 Monate
          \begin{itemize}
              \item
          \end{itemize}
\end{itemize}

\textbf{Geschäftsleitung OpenAI}
\begin{itemize}
    \item Statement an Geschäftsleitung mit Fakten (Kennzahl) 
    \begin{itemize}
        \item Alles was rechtmässig ist soll zulässig sein
        \item Wir haben Sicherheitsrichtlinien und Werte
    \end{itemize}
    \item konkrete Forderung an Geschäftsleitung
    \begin{itemize}
        \item keine
    \end{itemize}
    \item Was wäre das stärkste Argument der Gegenseite
    \begin{itemize}
        \item keine
    \end{itemize}
    \item Eine Erfolgskennzahl der letzten 12 Monate
    \begin{itemize}
        \item keine
    \end{itemize}
\end{itemize}
\textbf{Investorinnen und Verwaltungsrat} und
\textbf{Kundschaft und Belegschaft}
Nicht zum Sprechen gekommen. Debatte war nicht wirklich einleuchtend noch hilfreich... 

\pagebreak
\subsubsection{Scheinwiderspruch}

Trotz Absicht von 2.5 Mio Kunden ChatGPT Abo zu kündigen, wuchs die tägliche Nutzung um 0.3\%. Es wurde aber nicht berücksichtigt, dass das Unternehmen bereits in der vergangenen Periode stark gewachsen war.
Probleme: 
\begin{itemize}
    \item Parallele Trends sind nicht überprüfbar
\item Gemini parallel entwickelt — bleibt ungeprüft. Das ist die grösste Schwäche.
\item 3 OpenAI und Anthropic sind beide behandelt
\begin{itemize}
\item Nur in entgegengesetzter Richtung. Sie taugen deshalb nicht als Kontrolle füreinander
\end{itemize}
    \item Zeitfenster der Messung Januar-Februar beinhaltete auch andere Ereignisse die den Werten beeinflusste
\end{itemize}


\subsubsection{6 Risiken einer Wertepositionierung}

\subsubsection{4 Präventionspunkte einer Positionierung}
\begin{itemize}
    \item Kodifizieren, bevor es teuer wird (Werte und Konsequenzen abwägen)
    \item Eng und überprüfbar formulieren
    \begin{itemize}
        \item Anthropic lehnte nicht militärische Nutzung allgemein sondern 2 spezifische Anwendungen
    \end{itemize}
    \item Den Preis vorab beziffern
    \item Den Gegenpol nicht dem Gegner überlassen (mit beiden Seiten der Werte-Karte kommunizieren)
\end{itemize}

Vorlesung ging nur bis Folie 38 (wurde wegen Debatte verkürzt)